\chapter{Технологический раздел}
\section{Средства реализации}
Алгоритмы для данной лабораторной работы были реализованы на языке C++ при использовании компилятора gcc версии 13.3.1, так как в стандартной библиотеке приведённого языка присутствуют все необходимые структуры данных для реализации алгоритмов линейного и бинарного поиска.

\section{Реализация алгоритмов}
Листинги исходного кода функций линейного и бинарного поиска представлены в приложении~\ref{a:appendices}.

\section{Тестирование}

В таблице~\ref{t:mod_tests} приведены модульные тесты для разработанных алгоритмов поиска. Тестирование проводилось с помощью gtest, листинг исходного кода модульных тестов приведён в приложении~\ref{a:appendices}. Все тесты пройдены успешно.

\begin{table}[!hb]
	\centering
	\small
	\begin{threeparttable}
		\caption{Модульные тесты}
		\label{t:mod_tests}
		\begin{tabular}{|c|c|c|c|c|}
			\hline
			\bfseries Исходный массив
			& \bfseries X
			& \parbox[c][1.5cm][c]{2.5cm}{\centering \bfseries Ожидаемый \\ \bfseries результат}
			& \multicolumn{2}{c|}{\bfseries Полученный результат} \\ \cline{4-5}
			& & 
			& \parbox[c][1.5cm][c]{2.5cm}{\centering \bfseries Полный \\ \bfseries перебор}
			& \parbox[c][1.5cm][c]{2.5cm}{\centering \bfseries Бинарный \\ \bfseries поиск} \\
			\hline
			1 & 1 & 0 & 0 & 0 \\
			\hline
			123 & 1 & -1 & -1 & -1 \\
			\hline
			-999 -98 -34 0 23 45 67 99 111 999 & -999 & 0 & 0 & 0 \\
			\hline
			-999 -98 -34 0 23 45 67 99 111 999 & 45 & 5  & 5 & 5 \\
			\hline
			-999 -98 -34 0 23 45 67 99 111 999 & 999 & 9  & 9 & 9 \\
			\hline
			-999 -98 -34 0 23 45 67 99 111 999 & 1 & -1 & -1 & -1 \\
			\hline
		\end{tabular}  
	\end{threeparttable}
\end{table}
